\documentclass{standalone}
\usepackage[european]{circuitikz}

\begin{document}
	\begin{circuitikz} 
		\draw
		(0,0) to[V, l_={$U_0=10V$}] ++(0, -2)
		to[short] ++(0,-1)
		node[rground]{}
		(0,0) to[short] ++(3,0)
		to[pR, l_={$R_1=10k\Omega$}] ++(0, -2)
		to[short] ++(0,-1)
		node[rground]{}

		++(.5, 2) to[short] ++(1, 0) coordinate(a)
		to[R, l={$R_2=1k\Omega$}] ++(2, 0) coordinate(b)
		to[short] ++(1,0) coordinate(c)
		to[D, l=D] ++(0, -2) coordinate(d)
		node[rground]{}
		;


		\draw[draw=red, line width=1pt] 
		(a) to[short, *-] ++(0, 2) coordinate (K1a)
		to[voltmeter, color=red] (K1a-|b) coordinate (K1b)
		node[above, color=red, yshift=.4cm] at ($ (K1a)!0.5!(K1b) $) {Kanal 1}
		(K1b) to[short, -*] (b)
		;

		\draw[draw=blue, line width=1pt] 
		(c) to[short, *-] ++(2, 0) coordinate (K2a)
		to[voltmeter, color=blue] (K2a|-d) coordinate (K2b)
		node[above, color=blue, xshift=.4cm, rotate=-90] at ($ (K2a)!0.5!(K2b) $) {Kanal 2}
		(K2b) to[short, -*] (d)
		;
		
	\end{circuitikz}
	
	
\end{document}
